\section{F-SM-44: automatic low-degree compatibility and the precise higher obstruction}

\subsection{Statement}
Let \(k\) be a commutative ring and \(\mathcal M\) a complete
and cocomplete \(k\)-linear abelian category with enough
projectives.  Suppose there are \(k\)-linear adjunctions
\[
 \varphi:\mathcal M\rightleftarrows\operatorname{Mod}_k:\kappa,
       \qquad \varphi\dashv\kappa\dashv\psi,
       \qquad \kappa\text{ exact}.                           \tag{SM44.1}
\]
For every set-indexed family \((X_s)_{s\in S}\), the canonical
maps
\[
 b_n:\bigoplus_{s\in S}L_n\varphi(X_s)
           \longrightarrow L_n\varphi\!\left(\bigoplus_{s\in S}X_s\right)
                                                               \tag{SM44.2}
\]
are monomorphisms for all \(n\ge0\), and are isomorphisms for
\(n=0,1\).  No AB4, tensor exactness, closedness, or interval
assumption is needed for this assertion.

More precisely, fix an injective cogenerator \(I\) of
\(\operatorname{Mod}_k\), put \(Y=\kappa I\), and let
\(C_n=\operatorname{coker}b_n\).  Restriction to the summands
gives a natural exact sequence
\[
\begin{split}
 0\longrightarrow\operatorname{Hom}_k(C_n,I)
 &\longrightarrow
 \operatorname{Ext}^n_{\mathcal M}\!\left(\bigoplus_sX_s,Y\right)\\
 &\longrightarrow\prod_s\operatorname{Ext}^n_{\mathcal M}(X_s,Y)
       \longrightarrow0 .
\end{split}                                                    \tag{SM44.3}
\]
For \(n=0,1\), the left term vanishes.  Thus the remaining
derived-coproduct condition is precisely the vanishing, in
degrees \(n\ge2\), of the classes in the middle group whose
restriction to every summand is zero.

If in addition the canonical unit counit
\[
                      L\varphi(\kappa k)\longrightarrow k[0]
                                                               \tag{SM44.4}
\]
is an equivalence, then for every scalar module \(U\) the
canonical maps satisfy
\[
         \varphi\kappa U\simeq U,
                    \qquad L_1\varphi(\kappa U)=0.            \tag{SM44.5}
\]
In particular, any failure of full scalar normalization under
the original assumptions begins in degree at least two.
The exact scalar embedding also induces
\[
 \operatorname{Ext}^1_k(U,V)
       \simeq\operatorname{Ext}^1_{\mathcal M}(\kappa U,\kappa V)
                                                               \tag{SM44.6}
\]
for all \(U,V\).

\subsection{Complete proof}
\paragraph{1. Products are exact, and scalar duality computes the relevant Ext.}
Products in \(\mathcal M\) preserve kernels.  To see that they
preserve epimorphisms, take a family of epimorphisms
\(A_s\to B_s\), and choose a projective epimorphism
\(P\to\prod_sB_s\).  Its component maps lift to \(A_s\)
because \(P\) is projective.  Their product gives a map
\(P\to\prod_sA_s\) whose composite to \(\prod_sB_s\)
is epic.  Hence \(\prod_sA_s\to\prod_sB_s\) is epic.
This proves exactness of products without an AB4 assumption.

The functor \(\varphi\) takes projective objects to projective
scalar modules: \(\operatorname{Hom}_k(\varphi P,-)
=\operatorname{Hom}_{\mathcal M}(P,\kappa(-))\) is exact.
For a projective resolution \(P_\bullet\to X\), the
adjunction and injectivity of \(I\) therefore give, naturally,
\[
 \operatorname{Ext}^n_{\mathcal M}(X,\kappa I)
   =H^n\operatorname{Hom}_k(\varphi P_\bullet,I)
   =\operatorname{Hom}_k(L_n\varphi X,I).                     \tag{SM44.7}
\]
An injective cogenerator always exists in the scalar category;
one choice is \(I=\operatorname{Hom}_{\mathbf Z}(k,\mathbf Q/\mathbf Z)\).
The adjunction with the underlying abelian group identifies its
Hom functor with \(\operatorname{Hom}_{\mathbf Z}(-,\mathbf Q/\mathbf Z)\),
which is exact and detects nonzero abelian groups.

\paragraph{2. Every family of extension classes is the restriction of one class.}
We prove surjectivity of
\[
 R_n:\operatorname{Ext}^n_{\mathcal M}\!\left(\bigoplus_sX_s,Y\right)
              \longrightarrow\prod_s\operatorname{Ext}^n_{\mathcal M}(X_s,Y)
                                                               \tag{SM44.8}
\]
for \(n\ge1\).  The empty family is immediate, so assume
\(S\ne\varnothing\).  Represent prescribed classes \(e_s\)
by exact Yoneda \(n\)-extensions
\[
 0\longrightarrow Y\longrightarrow E_{s,n-1}
      \longrightarrow\cdots\longrightarrow E_{s,0}
      \longrightarrow X_s\longrightarrow0.
                                                               \tag{SM44.9}
\]
Their termwise product is exact by Step~1 and represents a class
\[
             e_{\Pi}\in
 \operatorname{Ext}^n_{\mathcal M}\!\left(\prod_sX_s,\prod_sY\right).
\]
Since \(\kappa\) is a right adjoint, \(\prod_sY=\kappa(I^S)\).
In the scalar module category the canonical inclusion
\(I^{(S)}\hookrightarrow I^S\) is monic.  The finite-support
codiagonal \(I^{(S)}\to I\) extends by injectivity of \(I\)
to a map
\[
                   t:I^S\longrightarrow I,
                   \qquad t\iota_s=1_I\quad(s\in S),         \tag{SM44.10}
\]
where \(\iota_s\) is the coordinate inclusion.
Let \(j:\bigoplus_sX_s\to\prod_sX_s\) be the canonical
map.  It need not be monic, and no such assertion is used.
Pull back the right end of \(e_{\Pi}\) along \(j\), and
push out its left end along \(\kappa t\).  These endpoint
operations on Yoneda extensions give
\[
                   e=(\kappa t)_*\,j^*e_{\Pi}
              \in\operatorname{Ext}^n_{\mathcal M}
                         \!\left(\bigoplus_sX_s,Y\right).     \tag{SM44.11}
\]

We check each restriction, including its coefficient map.
Denote the coproduct inclusion by \(\mu_s\), the coordinate
inclusion into \(\prod_sX_s\) by \(v_s\), and the coordinate
inclusion \(Y\to\prod_sY\) by \(w_s=\kappa\iota_s\).
There is a morphism from the \(s\)-th extension (SM44.9) to
its product extension, using coordinate inclusions in every
term, with endpoint maps \(w_s,v_s\).  Functoriality of
Yoneda pushout and pullback gives
\[
                 v_s^*e_{\Pi}=(w_s)_*e_s.
\]
Because \(j\mu_s=v_s\) and \((\kappa t)w_s=1_Y\),
\[
 \mu_s^* e
   =(\kappa t)_*v_s^*e_{\Pi}
   =((\kappa t)w_s)_*e_s=e_s.                               \tag{SM44.12}
\]
This proves the claimed surjectivity in every degree.  The
choice of \(t\) need not be natural and does not affect the
surjectivity assertion; \(R_n\) itself is canonical.

\paragraph{3. In degree one, restriction is also injective.}
This step holds in every abelian category with the relevant
coproduct, for every coefficient object \(Y\).  An extension
\[
                0\longrightarrow Y\longrightarrow E
                   \xrightarrow{q}\bigoplus_sX_s\longrightarrow0
\]
whose pullback to each \(X_s\) splits admits maps
\(u_s:X_s\to E\) with \(qu_s=\mu_s\).  By the coproduct
universal property, these assemble to \(u:\bigoplus_sX_s\to E\)
with \(qu=1\).  The extension splits, proving injectivity
of \(R_1\).  In degree zero the ordinary Hom coproduct
identity is already an isomorphism.

\paragraph{4. Dualize the restriction maps.}
Under (SM44.7) and
\(\operatorname{Hom}_k(\bigoplus_sL_n\varphi X_s,I)
=\prod_s\operatorname{Hom}_k(L_n\varphi X_s,I)\),
the map \(R_n\) is precisely \(\operatorname{Hom}_k(b_n,I)\).
The exact faithful contravariant functor
\(\operatorname{Hom}_k(-,I)\) detects monomorphisms by
surjectivity, and detects isomorphisms.  Step~2 therefore
makes every \(b_n\) monic, and Step~3 makes \(b_1\) an
isomorphism.  The case \(n=0\) also follows directly from
preservation of coproducts by the left adjoint \(\varphi\).
Dualizing
\(0\to\bigoplus_sL_n\varphi X_s\xrightarrow{b_n}
L_n\varphi(\bigoplus_sX_s)\to C_n\to0\)
gives exactly (SM44.3).

Equivalently, for the canonical derived comparison
\[
 b:\bigoplus_s L\varphi(X_s)
                \longrightarrow L\varphi\!\left(\bigoplus_sX_s\right),
                                                               \tag{SM44.13}
\]
the homology of its cone in degree \(n\) is \(C_n\): the
long exact sequence of the cone and the injectivity of
\(b_n,b_{n-1}\) give this assertion.  Thus this cone has
zero homology in degrees zero and one, and full derived
coproduct compatibility is equivalent to \(C_n=0\) for
all \(n\ge2\).  Ordinary scalar-module coproducts are exact,
so their homology commutes with the coproduct in (SM44.13);
no derived coproduct assertion in \(\mathcal M\) is used here.

\paragraph{5. Canonical scalar normalization through degree one.}
Assume (SM44.4).  Its degree-zero map says that the ordinary
counit \(\varphi\kappa k\to k\) is an isomorphism.
Both \(\varphi\) and \(\kappa\) preserve coproducts and
cokernels, since both are left adjoints in (SM44.1).
The same canonical counit is therefore an isomorphism on
free modules; applying it to a free presentation and taking
cokernels proves
\(\varphi\kappa U\xrightarrow{\sim}U\) for every \(U\).
In particular \(\kappa\), as the right adjoint of
\(\varphi\), is fully faithful.

For \(F=k^{(S)}\), one has
\(\kappa F=\bigoplus_s\kappa k\).  Applying the isomorphism
\(b_1\) and (SM44.4) gives
\(L_1\varphi(\kappa F)=0\).  For any scalar \(U\), choose
an exact sequence \(0\to K\to F\to U\to0\) with \(F\)
free.  Exactness of \(\kappa\) and the derived long exact
sequence give
\[
 0=L_1\varphi(\kappa F)\longrightarrow L_1\varphi(\kappa U)
       \longrightarrow\varphi\kappa K\longrightarrow\varphi\kappa F.
\]
Under the canonical counit identifications, the last map is
the monomorphism \(K\to F\).  Thus
\(L_1\varphi(\kappa U)=0\), proving (SM44.5).

Finally, given an extension
\(0\to\kappa V\to Z\to\kappa U\to0\), the vanishing
just proved makes its image under \(\varphi\) the exact
sequence \(0\to V\to\varphi Z\to U\to0\).
The unit \(Z\to\kappa\varphi Z\) gives a morphism from
the original extension to the image of this scalar extension.
Its endpoint maps are isomorphisms by the adjunction identities,
so its middle map is an isomorphism.  Applying \(\varphi\)
and \(\kappa\) to extensions therefore gives inverse maps
on their classes, including the Baer operations, which proves
(SM44.6).  \(\square\)

\subsection{Sharp scope and the remaining original step}
The tensor and interval hypotheses have not been used to prove
this result.  Thus it applies in particular to the original
abstract setup.  It reduces the unresolved coproduct question
to \emph{surjectivity} of (SM44.2) in degrees \(n\ge2\),
or equivalently to injectivity of the corresponding higher
restriction maps (SM44.8).  It does not prove that the tensor
or interval kills those higher kernels.

The bound cannot be improved by the additive assumptions alone:
in F-SM-34, for the family of countably many copies of
\(\kappa k\), the degree-two map is
\(0\to E(k^{(\mathbb N)})\ne0\).  This is a precise additive
sharpness example, not a full symmetric-tensor counterexample.

\subsection{Primary reference and check scope}
Argud\'in-Monroy, A. (2022). The Yoneda Ext and arbitrary
coproducts in abelian categories. \emph{Glasgow Mathematical
Journal, 64}(2), 277--291.
\url{https://doi.org/10.1017/S0017089521000070}.
Version inspected: arXiv:1904.12182v2, September~8, 2020,
\url{https://arxiv.org/pdf/1904.12182}.
Proposition~3.11(b), p.~12, proves the degree-one injection by
assembling sections; the discussion after Theorem~3.13, p.~14,
explicitly states its validity for AB3 and separates the
higher-degree problem.  Only these precise source passages
and the publisher's bibliographic metadata were inspected for
this lemma.  The product-extension construction and its
application to the two scalar adjoints are proved in full above;
no general AB4 theorem is being applied to \(\mathcal M\).
